%\documentclass[11pt,a4paper,uplatex,dvipdfmx]{ujarticle} 		% for uplatex
\documentclass[11pt,a4j,dvipdfmx]{jarticle} 					% for platex
\input{pieces/form00_header} % pieces
\input{pieces/kakenhi7} % pieces
\input{pieces/form01_header} % pieces
\input{pieces/hook3} % pieces
%#Name: kokusai_renkei
\input{pieces/form04_jsps_headers} % pieces
\input{pieces/form04_kokusai_renkei_header} % pieces
% ===== Global definitions for the Kakenhi form ======================
%　基本情報
%
%------ 研究課題名  -------------------------------------------
\newcommand{\研究課題名}{象の卵}

%----- 研究機関名と研究代表者の氏名-----------------------
\newcommand{\研究機関名}{逢坂大学}
\newcommand{\研究代表者氏名}{湯川秀樹}
\newcommand{\me}{\underline{\underline{H.~Yukawa}}} 
\input{pieces/inst_general_images} % pieces
% user07_header
% ===== my favorite packages ====================================
% ここに、自分の使いたいパッケージを宣言して下さい。
\usepackage{wrapfig}
%\usepackage{amssymb}
%\usepackage{mb}
%\DeclareGraphicsRule{.tif}{png}{.png}{`convert #1 `dirname #1`/`basename #1 .tif`.png}
\usepackage{lineno}

% ===== my personal definitions ==================================
% ここに、自分のよく使う記号などを定義して下さい。
\newcommand{\klpionn}{K_L \to \pi^0 \nu \overline{\nu}}
\newcommand{\kppipnn}{K^+ \to \pi^+ \nu \overline{\nu}}

% ----- 業績リスト用 -------------
\newcommand{\paper}[6]{%
	% paper{title}{authors}{journal}{vol}{pages}{year}
	\item ``#1'', #2, #3 {\bf #4}, #5 (#6).			% お好みに合わせて変えてください。
}

\newcommand{\etal}{\textit{et al.\ }}
\newcommand{\ca}[1]{*#1}	% corresponding author;   \ca{\yukawa}  みたいにして使う
\newcommand{\invitedtalk}{招待講演}

\newcommand{\yukawa}{H.~Yukawa}					% no underline
%\newcommand{\yukawa}{\underline{\underline{H.~Yukawa}}}	% with 2 underlines
\newcommand{\tomonaga}{S.~Tomonaga}

\newcommand{\prl}{Phys.\ Rev.\ Lett.\ }		% よく使う雑誌も定義すると楽

% ===== 欄外メモ ==================
\newcommand{\memo}[1]{\marginpar{#1}}
%\renewcommand{\memo}[1]{}	% 全てのメモを表示させないようにするには、行頭の"%"を消す

%\input{../../sample/simple/contents}	% skip
\input{pieces/hook5} % pieces

\begin{document}
\input{pieces/hook7} % pieces
%#Split: importance_need  
%#PieceName: pimportance_need
\input{pieces/pimportance_need_00}
\section{１　国際共同研究の概要及び意義・必要性など}
%　　　　＜＜最大　3ページ＞＞

%s08_need_KKD_with_abstract
%begin 作成留意事項 ＝＝＝＝＝＝＝＝＝＝＝＝＝＝＝＝＝＝＝＝
\JSPSInstructions	% <-- 留意事項。これは消すか、コメントアウトしてください。
%end 作成留意事項 ＝＝＝＝＝＝＝＝＝＝＝＝＝＝＝＝＝＝＝＝
\noindent
\textbf{（概要）}\\
%begin 国際的研究の意義と必要性の概要 ＝＝＝＝＝＝＝＝＝＝＝＝＝＝＝＝＝＝＝＝
	象の生息地に行かないと真の研究はできない。
	象牙の価値を下げないと、象の生息地から象が消える。
	象牙の価値を下げる最良の方法は、国際協力で象の卵の殻から象牙の代替品を作ることである。
%end 国際的研究の意義と必要性の概要 ＝＝＝＝＝＝＝＝＝＝＝＝＝＝＝＝＝＝＝＝
\\

\noindent
\textbf{（本文）}\\
%begin 国際的研究の意義と必要性 ＝＝＝＝＝＝＝＝＝＝＝＝＝＝＝＝＝＝＝＝
	国際的研究の意義と必要性は．．．
%end 国際的研究の意義と必要性 ＝＝＝＝＝＝＝＝＝＝＝＝＝＝＝＝＝＝＝＝

\input{pieces/pimportance_need_01}

%#Split: purpose_plan  
%#PieceName: ppurpose_plan
\input{pieces/ppurpose_plan_00}
\section{２　国際共同研究の研究目的、研究方法など}
%　　　　＜＜最大　3ページ＞＞

%s02_purpose_plan_KKD_no_abstract
%begin 研究目的と研究計画留意事項なし ＝＝＝＝＝＝＝＝＝＝＝＝＝＝＝＝＝＝＝＝
	 象の卵の研究目的は．．．

	象の卵の研究計画は．．．

	\vspace{1cm}
	\begin{thebibliography}{99}
		\bibitem{teramura} 寺村輝夫、「ぼくは王様 - ぞうのたまごのたまごやき」.
	\end{thebibliography}
%end 研究目的と研究計画留意事項なし ＝＝＝＝＝＝＝＝＝＝＝＝＝＝＝＝＝＝＝＝

\input{pieces/ppurpose_plan_01}

%#Split: preparation  
%#PieceName: ppreparation
\input{pieces/ppreparation_00}
\section{３　海外共同研究者の役割及び準備状況}
%　　　　＜＜最大　2ページ＞＞

%s04_prep_KKD
%begin 海外共同研究者の役割と準備状況 ＝＝＝＝＝＝＝＝＝＝＝＝＝＝＝＝＝＝＝＝
	先進国の粉体圧縮加工と高分子化学の研究者は、象の卵の殻から象牙と同等な材質を作る技術を開発する。
	象の生息地の生物学者は、密猟ビジネスに詳しい社会学者と協力して、
	新しい象の卵のビジネスモデルを構築する。
	
	海外共同研究者とは、毎日つぶやきあっている。

	J. Gutenberg, ``Forming powder into solid'', J. Powder \textbf{800}, 999 (2016).
%end 海外共同研究者の役割と準備状況 ＝＝＝＝＝＝＝＝＝＝＝＝＝＝＝＝＝＝＝＝

\input{pieces/ppreparation_01}

%#Split: abilities  
%#PieceName: pabilities
\input{pieces/pabilities_00}
\section{４　応募者の研究遂行能力及び研究環境}
%　　　　＜＜最大　2ページ＞＞

% s14_abilities
%begin 業績留意事項 ＝＝＝＝＝＝＝＝＝＝＝＝＝＝＝＝＝＝＝＝
\PapersInstructions		% <-- 留意事項。これは消すか、コメントアウトしてください。
%end 業績留意事項 ＝＝＝＝＝＝＝＝＝＝＝＝＝＝＝＝＝＝＝＝
%begin 応募者の研究遂行能力及び研究環境 ＝＝＝＝＝＝＝＝＝＝＝＝＝＝＝＝＝＝＝＝
	応募者は過去20年間、７つの海を隅から隅まで航海し、
	浅瀬から深海まで潜り、文字通り東西南北上下の３次元で
	シロナガスクジラの卵の探索を行ってきた(業績\ref{pub:whale})。
	シロナガスクジラに飲み込まれそうになったり、海賊に捕まるなどの危険な目にも
	あったが、それにもめげず、研究を遂行してきた強靭な能力を有する。

	シロナガスクジラの卵を探すために用いていたソナーと双眼鏡、及び
	シロナガスクジラの卵を引き上げるために用意していた大きな網は、
	そのまま使える。
%end 応募者の研究遂行能力及び研究環境 ＝＝＝＝＝＝＝＝＝＝＝＝＝＝＝＝＝＝＝＝
%begin 研究業績リスト ＝＝＝＝＝＝＝＝＝＝＝＝＝＝＝＝＝＝＝＝
	\begin{enumerate}
		\paper{Search for whale eggs}{\yukawa\ \etal}{Rev.\ Oceanic Mysteries}{888}{99}{2017}
			\label{pub:whale}
				
		\paper{Theory of Elephant Eggs}{\yukawa, Kara Juzo \etal}{\prl}{800}{800-804}{2005}
			\label{pub:theoegg}
				
		\paper{仔象は死んだ}{Kobo Abe}{安部公房全集}{26}{100-200}{2004}
		
		\paper{The Elephant's Child (象の鼻はなぜ長い)}{R.~Kipling}{Nature}{999}{777-799}{2003}

		\paper{You can't Lay an Egg If You're an Elephant}{F.~Ehrlich}
			{JofUR\\({\tt www.universalrejection.org})}{{\bf N/A}}{2002}
		
		% 下のように書いてもいいけど、めんどくさいし、表示の仕方を変えようとしたら大変。
		\item ``Egg of Elephant-Bird'', 
				\underline{A.~Cooper},
				Nature, {\bf 409}, 704-707 (2001).	% 	
		\input{jack_pub}	% << only for demonstration. Please delete it or comment it out.
	\end{enumerate}
%end 研究業績リスト ＝＝＝＝＝＝＝＝＝＝＝＝＝＝＝＝＝＝＝＝

\input{pieces/pabilities_01}

%#Split: rights  
%#PieceName: prights
\input{pieces/prights_00}
\section{５　人権の保護及び法令等の遵守への対応 }
%　　　　＜＜最大　1ページ＞＞

% s09_rights
%begin 人権の保護及び法令等の遵守への対応 ＝＝＝＝＝＝＝＝＝＝＝＝＝＝＝＝＝＝＝＝
	象の卵のES細胞の培養、象のクローンの生成などは行わない。

	\LaTeX の便利な機能については、\texttt{egg\_***.tex} や\texttt{sample\_pdf/egg\_***.pdf}をご覧ください。
%end 人権の保護及び法令等の遵守への対応 ＝＝＝＝＝＝＝＝＝＝＝＝＝＝＝＝＝＝＝＝

\input{pieces/prights_01}

%#Split: 99_tail
\input{pieces/hook9} % pieces
\end{document}

